\documentclass[11pt]{article}

\usepackage[margin=1in]{geometry}
\usepackage{amsmath, amssymb, amsthm}
\usepackage{mathtools}
\usepackage[colorlinks=true, linkcolor=blue, urlcolor=blue, citecolor=blue]{hyperref}

\title{\textbf{Inequality constraints in QUBO}}
\author{Sanjay Suresh}
\date{April 2026}

\begin{document}
\maketitle

To map an inequality like $Ax \le b$ onto a QUBO for an annealer, introduce binary-encoded slack
variables and fold the constraint into the objective as a quadratic penalty term. The equality
$Ax + s = b$ becomes
\[ (Ax + s - b)^2, \]
which is zero exactly when the constraint is satisfied.

The penalty weight matters more than people expect. Too small and the global optimum quietly
violates the constraint; too large and it dominates the energy landscape, flattening out the
differences between good feasible solutions so the annealer cannot tell them apart. There is a
usable window between ``constraint is respected'' and ``objective still has resolution'' --- sweep
the weight to find it rather than guessing once.

Slack encoding also burns qubits fast, so prefer the tightest bound you can prove on the slack range
before you discretize it.

\end{document}
